\section{Experiments}
\label{sec:experiments}

\subsection{Main Results}
\label{sec:main-results}

\Cref{tab:main-results} compares \method{} with three baselines on four benchmarks, and \cref{fig:benchmark-scores} plots the same scores against the strongest baseline.
A results paragraph opens with its finding rather than with the table it comes from.

\begin{table*}[t]
  \centering
  \caption{Main results on four benchmarks. The best score in each column is in bold and the second best is underlined. All numbers are placeholders.}
  \label{tab:main-results}
  % Main results; the caption and label are in sections/experiments.tex.
\small
\begin{tabular}{@{}lccccc@{}}
  \toprule
  Method & Bench-A & Bench-B & Bench-C & Bench-D & Avg. \\
  \midrule
  Base model         & 41.2 & 55.0 & 30.4 & 62.1 & 47.2 \\
  Prompting baseline & 44.8 & 57.3 & 33.9 & 63.0 & 49.8 \\
  Published method   & \second{49.5} & \second{60.2} & \second{37.1} & \second{66.4} & \second{53.3} \\
  \rowcolor{labtint}
  \method{}          & \best{53.1} & \best{63.7} & \best{41.8} & \best{68.0} & \best{56.7} \\
  \bottomrule
\end{tabular}

\end{table*}

\begin{figure}[t]
  \centering
  % Scores of the method and the strongest baseline, from tables/main-results.tex.
\begin{tikzpicture}
  \begin{axis}[
      width=\linewidth, height=4.4cm, ybar=2pt, bar width=10pt,
      ymin=0, ymax=80, ytick={0,20,40,60,80}, enlarge x limits=0.12,
      /pgf/number format/assume math mode=true,
      symbolic x coords={Bench-A,Bench-B,Bench-C,Bench-D,Average}, xtick=data,
      axis lines*=left, axis line style={draw=labrule}, tick style={draw=none},
      ymajorgrids, grid style={draw=black!8},
      ticklabel style={font=\sffamily\scriptsize, text=labgrey},
      ylabel={Score}, ylabel style={font=\sffamily\scriptsize, text=labgrey},
      legend style={font=\sffamily\scriptsize, draw=none, fill=none,
        at={(0.5,1.02)}, anchor=south, legend columns=2,
        /tikz/every even column/.append style={column sep=12pt}},
      area legend,
      nodes near coords style={font=\sffamily\tiny, text=labink},
      every node near coord/.append style={
        /pgf/number format/.cd, fixed, fixed zerofill, precision=1, assume math mode=true}]
    \addplot[fill=chartblue, draw=none]
      coordinates {(Bench-A,49.5) (Bench-B,60.2) (Bench-C,37.1) (Bench-D,66.4) (Average,53.3)};
    \addplot[fill=chartpurple, draw=none, nodes near coords]
      coordinates {(Bench-A,53.1) (Bench-B,63.7) (Bench-C,41.8) (Bench-D,68.0) (Average,56.7)};
    \legend{Strongest published baseline, \method{}}
  \end{axis}
\end{tikzpicture}

  \caption{Scores of \method{} and the strongest published baseline, from the placeholder numbers of \cref{tab:main-results}. The chart uses the class colours \texttt{chartpurple} and \texttt{chartblue}.}
  \label{fig:benchmark-scores}
\end{figure}

\begin{takeaway}[Takeaway]
State the finding in one or two sentences with the number that supports it, for example that \method{} improves the average score by 3.4 points over the strongest published baseline in \cref{tab:main-results}.
\end{takeaway}

\subsection{Prompts}
\label{sec:prompts}

The \texttt{promptbox} environment sets a prompt or a model output in Inconsolata, and verbatim content is allowed inside it.

\begin{promptbox}[Prompt template for the judge model]
You are grading a response to the task below.\\
Task: \{task\}\\
Response: \{response\}\\
Reply with a score from 1 to 10 and one sentence of justification.
\end{promptbox}
